% This is text associated with the open textbook "Open Electromagnetics"
% See immediately below for author, date
% <License info goes here (boilerplate, not read by project compiler)>

%ORB TITLE Permittivity    
%ORB AUTHOR 0       # S.W. Ellingson, Virginia Tech (ellingson@vt.edu)
%ORB DATE   20180731
%ORB ABSTRACT Introduction to permittivity, using concept of relationship between electric field intensity and charge

\label{m0008_Permittivity} % <-- Must match name of this file and module directory name.  Don't mess with this!
\index{permittivity|(}

Permittivity describes the effect of material in determining the electric field in response to electric charge.  For example, one can observe from laboratory experiments %\footnote{Also worth noting that this equation also follows directly from Maxwell's Equations, described in the next chapter.} 
that a particle having charge $q$ gives rise to the electric field 
%
\begin{equation}
{\bf E} = \hat{\bf R} ~ q ~ \frac{1}{4\pi R^2} ~ \frac{1}{\epsilon}
\end{equation}
%
where $R$ is distance from the charge, 
$\hat{\bf R}$ is a unit vector pointing away from the charge,
and $\epsilon$ is a constant that depends on the material.
Note that ${\bf E}$ increases with $q$, which makes sense since electric charge is the source of ${\bf E}$.  
Also note that ${\bf E}$ is inversely proportional to $4\pi R^2$, indicating that ${\bf E}$ decreases in proportion to the area of a sphere surrounding the charge -- a principle commonly known as the \emph{inverse square law}.
\index{inverse square law}
The remaining factor $1/\epsilon$ is the constant of proportionality, which captures the effect of material.   
Given units of V/m for ${\bf E}$ and C for $Q$, we find that $\epsilon$ must have units of farads per meter (F/m). 
(To see this, note that 1~F $=$ 1~C/V.)    

%ORB HIGHLIGHT
\begin{mdframed}[backgroundcolor=yellow!20,nobreak=true] 
Permittivity ($\epsilon$, F/m) describes the effect of material in determining the electric field intensity in response to charge. 
\end{mdframed}

In free space (that is, a perfect vacuum), we find that $\epsilon = \epsilon_0$ where:
%
\begin{equation}
\epsilon_0 \cong 8.854 \times 10^{-12} ~\mbox{F/m}
\end{equation}
% 
The permittivity of air is only slightly greater, and usually can be assumed to be equal to that of free space.
In most other materials, the permittivity is significantly greater; that is, the same charge results in a weaker electric field intensity.  

It is common practice to describe the permittivity of materials relative to the permittivity of free space.
This \emph{relative permittivity} 
\index{permittivity!relative}
\index{relative permittivity|see{permittivity, relative}}
is given by:
%
\begin{equation}
\epsilon_r \triangleq \frac{\epsilon}{\epsilon_0}
\end{equation}
%  
Values of $\epsilon_r$ for a few representative materials is given in Appendix~\ref{m0135_Table_of_Permittivities}.  
Note that $\epsilon_r$ ranges from 1 (corresponding to a perfect vacuum) to about 60 or so in common engineering applications.  
Also note that relative permittivity is sometimes referred to as \emph{dielectric constant}. 
\index{dielectric constant}
This term is a bit misleading, however, since permittivity is a meaningful concept for many materials that are not dielectrics.

%%%%%%%%%%%%%%%%%

{\bf Additional Reading:}
\begin{itemize}
\item ``\href{https://en.wikipedia.org/wiki/Permittivity}{Permittivity}'' on Wikipedia.
\item Appendix~\ref{m0135_Table_of_Permittivities} (``Permittivity of Some Common Materials'').
\item ``\href{https://en.wikipedia.org/wiki/Inverse-square_law}{Inverse square law}'' on Wikipedia.
\end{itemize}

\index{permittivity|)}


